%% =============================================================
%%  Informe de revisi\'on y hoja de ruta de investigaci\'on
%%  para el manuscrito RutAI (R1)
%%  Formato: art\'iculo autocontenido, compilable con pdflatex
%% =============================================================

\documentclass[11pt,a4paper]{article}

% -- Codificaci\'on y fuentes (primero siempre)
\usepackage[utf8]{inputenc}
\usepackage[T1]{fontenc}
\usepackage{lmodern}

% -- Idioma
\usepackage[spanish,es-tabla,es-nolists]{babel}

% -- M\'argenes y tipograf\'ia
\usepackage[margin=2.5cm]{geometry}
\usepackage{microtype}
\usepackage{parskip}
\setlength{\parindent}{0pt}
\setlength{\parskip}{0.5em}

% -- Tablas y listas
\usepackage{booktabs}
\usepackage{array}
\usepackage{longtable}
\usepackage{enumitem}
\setlist[itemize]{leftmargin=*, topsep=2pt, itemsep=2pt}
\setlist[enumerate]{leftmargin=*, topsep=2pt, itemsep=2pt}
\setlist[description]{leftmargin=0pt, labelindent=0pt}

% -- Cajas destacadas para veredicto/recomendaciones
\usepackage[most]{tcolorbox}
\tcbuselibrary{skins,breakable}

\newtcolorbox{verdictbox}[1][]{
	enhanced, breakable,
	colback=gray!6, colframe=black!70,
	boxrule=0.4pt, arc=2pt,
	left=8pt, right=8pt, top=6pt, bottom=6pt,
	fonttitle=\bfseries,
	title={#1}
}

\newtcolorbox{actionbox}[1][]{
	enhanced, breakable,
	colback=blue!3, colframe=blue!40!black,
	boxrule=0.4pt, arc=2pt,
	left=8pt, right=8pt, top=6pt, bottom=6pt,
	fonttitle=\bfseries\small,
	title={#1}
}

% -- Encabezados y pies
\usepackage{fancyhdr}
\pagestyle{fancy}
\fancyhf{}
\fancyhead[L]{\small Informe t\'ecnico de revisi\'on --- RutAI R1}
\fancyhead[R]{\small \thepage}
\fancyfoot[C]{\small Elaborado como revisor senior JCR Q1/Q2}
\renewcommand{\headrulewidth}{0.3pt}

% -- Hyperref al final
\usepackage[hidelinks, colorlinks=false]{hyperref}

% -- T\'itulos de secci\'on m\'as compactos
\usepackage{titlesec}
\titleformat*{\section}{\large\bfseries}
\titleformat*{\subsection}{\normalsize\bfseries}
\titleformat*{\subsubsection}{\normalsize\itshape}

% -- C\'odigo en l\'inea
\usepackage{listings}
\lstset{
	basicstyle=\ttfamily\footnotesize,
	breaklines=true,
	frame=single,
	framerule=0.2pt,
	rulecolor=\color{black!40},
	backgroundcolor=\color{gray!5},
	xleftmargin=0.5em,
	xrightmargin=0.5em,
	aboveskip=0.5em,
	belowskip=0.5em
}
\usepackage{xcolor}
\usepackage{xspace}

% =============================================================
\title{Informe de revisi\'on por pares \\ y hoja de ruta de investigaci\'on \\[6pt]
	\large Manuscrito: \emph{RutAI --- A Component-Based Mobile Framework}}
\author{Revisor senior \\ Revista objetivo: \emph{Pervasive and Mobile Computing} (Elsevier, JCR Q2)}
\date{Abril 2026}

\newcommand{\RUTAI}{RutAI\xspace}
\newcommand{\GPS}{GPS\xspace}
\newcommand{\IOT}{IoT\xspace}
\newcommand{\LBS}{LBS\xspace}
\newcommand{\POI}{POI\xspace}
\newcommand{\UCB}{UCB\xspace}
\newcommand{\MAB}{MAB\xspace}
\newcommand{\TAM}{TAM\xspace}
\newcommand{\CBSE}{CBSE\xspace}
\newcommand{\MVVM}{MVVM\xspace}


\begin{document}
	\maketitle
	
	\begin{verdictbox}[Recomendaci\'on global]
		\textbf{Major Revision} (segunda ronda). El manuscrito ha mejorado sustancialmente respecto a iteraciones previas, pero persisten brechas evidenciales serias que impiden una aceptaci\'on directa.
	\end{verdictbox}
	
	\tableofcontents
	\vspace{1em}
	
	% =============================================================
	\part*{PARTE I --- Informe de evaluador}
	\addcontentsline{toc}{section}{PARTE I --- Informe de evaluador}
	% =============================================================
	
	\section{Fortalezas reconocidas}
	
	Antes de las cr\'iticas es justo reconocer lo que el manuscrito hace bien:
	
	\begin{itemize}
		\item Estructura IMRaD extendida correcta y alineada con Elsevier CAS.
		\item Redacci\'on en espa\~nol de alta calidad, tono humanizado apropiado.
		\item Bibliograf\'ia verificable con DOIs reales (48 entradas).
		\item Honestidad intelectual al reportar HTMT$>0{,}85$ y reconocer fallo de validez discriminante.
		\item Reformulaci\'on del UCB con t\'ermino de entrop\'ia como aporte genuino y medible.
		\item Integraci\'on de tres m\'odulos (bandit + geofencing + WebSocket) efectivamente \'unica en la literatura revisada.
	\end{itemize}
	
	\section{Observaciones cr\'iticas bloqueantes}
	
	\subsection{El banco de pruebas t\'ecnico no es reproducible con la informaci\'on dada}
	
	La Secci\'on 3.2.7 promete un protocolo reproducible, pero la Tabla~6 reporta cifras sin:
	
	\begin{itemize}
		\item \textbf{Descripci\'on del ``perfil sint\'etico derivado de trazas reales anonimizadas''} --- ?`qu\'e trazas?, ?`de d\'onde?, ?`cu\'antos usuarios?, ?`cu\'antos destinos?, ?`qu\'e recompensa real se us\'o para simular la respuesta del usuario frente a cada brazo?
		\item \textbf{Funci\'on de recompensa del entorno simulado}: sin ella, el n\'umero de regret es irreproducible. Un revisor no puede verificar $153{,}6\pm 9{,}8$ sin saber contra qu\'e \'optimo se midi\'o.
		\item \textbf{Par\'ametros exactos de cada baseline}: prior de Thompson sampling (?`$\beta(1,1)$?, ?`Gaussiano?), semilla inicial, horizonte de la ventana de entrop\'ia, tama\~no de la malla para $\lambda$.
		\item \textbf{C\'odigo ni dataset liberados}: ``disponibles a solicitud razonable'' es inaceptable en PMC desde 2023. Se exige DOI p\'ublico (Zenodo/Figshare).
	\end{itemize}
	
	\subsection{Las mediciones de latencia y energ\'ia carecen de condiciones de contorno}
	
	La Secci\'on 4.1 reporta ``latencia p95 de REST 312\,ms'', pero no especifica:
	
	\begin{itemize}
		\item \textbf{Red de prueba} (Wi-Fi dom\'estica, 4G simulada, 5G, latencia local al servidor Render.com desde Ecuador/Espa\~na).
		\item \textbf{Distribuci\'on geogr\'afica} del cliente respecto al servidor: un servidor en Render.com US-East genera $\sim$150\,ms de RTT base desde Ecuador, lo que explica $\geq 50\%$ de los 312\,ms reportados.
		\item \textbf{Carga concurrente real} cuando se midi\'o (100 usuarios simult\'aneos en Locust, pero ?`tasa de arribo?, ?`patr\'on \emph{think-time}?).
		\item \textbf{\emph{Warm-up} del JIT y base de datos}: si los primeros 500 requests se incluyeron, el p95 est\'a contaminado.
	\end{itemize}
	
	\subsection{El consumo de bater\'ia es anecd\'otico, no medici\'on rigurosa}
	
	``Tres sesiones de dos horas'' con $n=1$ dispositivo (Pixel~7) \textbf{no es una medici\'on, es una observaci\'on}. Se requiere:
	
	\begin{itemize}
		\item \textbf{$N\geq 5$ dispositivos} cubriendo gama baja, media y alta (Android Go, \emph{midrange} y \emph{flagship}).
		\item \textbf{Control de variables}: brillo de pantalla fijo, modo avi\'on para aislar GPS, temperatura ambiente controlada.
		\item \textbf{Replicaci\'on suficiente} para reportar intervalo de confianza de energ\'ia; tres sesiones dan $\pm\sigma$ pero no un IC cre\'ible.
		\item \textbf{Comparaci\'on ortogonal}: no s\'olo continuo vs. adaptativo, sino tambi\'en frente a las APIs nativas de Android (FusedLocationProvider con \texttt{PRIORITY\_BALANCED\_POWER\_ACCURACY}), que es la baseline industrial real.
	\end{itemize}
	
	\subsection{La precisi\'on del geofencing es potencialmente enga\~nosa}
	
	La Tabla~6 reporta $94{,}1\%$ a 100\,m, $89{,}2\%$ a 50\,m y $71{,}4\%$ a 25\,m. Tres problemas:
	
	\begin{itemize}
		\item S\'olo se reporta \textbf{recall}, no \textbf{precisi\'on} ni \textbf{F1} ni \textbf{latencia del trigger} (tiempo desde el cruce real hasta la notificaci\'on). Para una app de seguridad, un trigger que llega 45\,s tarde es in\'util.
		\item \textbf{No hay desglose por velocidad}: $71{,}4\%$ a 25\,m y 8,3\,m/s significa que el dispositivo atraviesa la cerca en $\sim$3\,s: ?`cu\'antos fallos se deben a muestreo insuficiente vs. error de GPS?
		\item \textbf{150 cruces es muestra peque\~na} para un IC sobre una proporci\'on; la amplitud del IC de Wilson a $71{,}4\%$ con $n=50$ es $\pm 12{,}5$ puntos.
	\end{itemize}
	
	\subsection{El estudio TAM es d\'ebil y los autores lo reconocen, pero no ofrecen plan estad\'istico s\'olido de rescate}
	
	El reconocimiento honesto del HTMT$>0{,}85$ es apropiado, pero:
	
	\begin{itemize}
		\item Se promete replicaci\'on con $n\geq 100$ \textbf{como trabajo futuro}, no como parte de la revisi\'on. Un revisor estricto exigir\'a los datos ahora.
		\item \textbf{Sin validez discriminante}, las correlaciones reportadas ($r>0{,}94$) no son informativas: miden la misma cosa cinco veces.
		\item \textbf{El instrumento necesita reformulaci\'on}: los \'items por constructo probablemente solapan sem\'anticamente. Hay que refinar la operacionalizaci\'on de cada dimensi\'on, especialmente Seguridad Percibida vs. Utilidad Percibida.
	\end{itemize}
	
	\subsection{La afirmaci\'on de ``safer urban mobility'' en el t\'itulo no est\'a demostrada emp\'iricamente}
	
	El t\'itulo promete \emph{mobility m\'as segura}. Los resultados ofrecen:
	
	\begin{itemize}
		\item Percepci\'on de seguridad (TAM-PS) $\neq$ seguridad real.
		\item Entrop\'ia de rutas mayor $\neq$ reducci\'on medida de incidentes.
	\end{itemize}
	
	\textbf{Falta el experimento que cierre el bucle}: un estudio longitudinal con usuarios reales comparando tasa de victimizaci\'on, ansiedad auto-reportada, o al menos reducci\'on de predictibilidad en un uso sostenido de $\geq 4$ semanas.
	
	\section{Observaciones mayores}
	
	\subsection{Las entrevistas de requisitos (\texorpdfstring{$n=5$}{n=5}) siguen siendo fr\'agiles}
	
	Mencionar a Guest et al. (2006) sobre saturaci\'on tem\'atica no reemplaza alcanzar esa saturaci\'on. Se requieren \textbf{al menos 12 entrevistas} con muestreo te\'oricamente diverso (g\'enero, edad, modalidad de transporte, nivel de inseguridad percibida).
	
	\subsection{No hay an\'alisis de sensibilidad al par\'ametro \texorpdfstring{$\lambda$}{lambda} de la regularizaci\'on por entrop\'ia}
	
	Se fija $\lambda=0{,}25$ ``mediante b\'usqueda en malla sobre un conjunto piloto'', pero no se reporta la malla, el criterio de optimalidad ni una curva $\lambda$ vs. regret/entrop\'ia. Un revisor pedir\'a esa tabla porque $\lambda$ determina si el aporte existe.
	
	\subsection{Ausencia total de capturas de pantalla de la app real}
	
	Para una publicaci\'on en PMC sobre un sistema m\'ovil, \textbf{se espera ver el sistema}. El manuscrito s\'olo tiene tres diagramas TikZ vectoriales. Faltan:
	
	\begin{itemize}
		\item Captura de la pantalla de selecci\'on de ruta (con las tres alternativas UCB).
		\item Captura del mapa con geocerca activa.
		\item Captura del grupo colaborativo con pings en tiempo real.
	\end{itemize}
	
	\subsection{\texttt{\textbackslash printcredits} duplica informaci\'on y los ORCID son \emph{placeholder}}
	
	Los ORCID \texttt{0000-0000-0000-0000} no son v\'alidos para Editorial Manager de Elsevier y delatan que el env\'io no se ha preparado.
	
	\subsection{Faltan declaraciones obligatorias de Elsevier 2024}
	
	\begin{itemize}
		\item Declaraci\'on de uso de IA generativa (pol\'itica Elsevier obligatoria desde 2023).
		\item Declaraci\'on CRediT completa (aunque los autores ya asignaron roles, falta formalizar \emph{``All authors read and approved the final manuscript''}).
		\item N\'umero de aprobaci\'on del Comit\'e de \'Etica para el estudio con $n=25$ (ausente).
	\end{itemize}
	
	\subsection{La Figura 4 mezcla lo num\'erico con lo textual de manera sub\'optima}
	
	Las etiquetas $4{,}14$, $4{,}18$, etc. sobre las barras son innecesarias cuando ya est\'an en la Tabla~7. Mejor: un \emph{violin plot} o \emph{dot plot} que muestre los 25 puntos individuales con \emph{overlay} del IC.
	
	\section{Observaciones menores}
	
	\begin{itemize}
		\item \textbf{Veredicto del idioma}: el manuscrito oscila entre espa\~nol para el cuerpo e ingl\'es para las figuras y tablas. Esto es inusual. Si se env\'ia a PMC (anglo), debe estar \textbf{todo en ingl\'es}; si se env\'ia a una revista hispanohablante, todo en espa\~nol.
		\item ``\emph{Una observaci\'on leg\'itima del revisor fue\ldots}'' en la Secci\'on 3.2.3 rompe la cuarta pared. Eliminar; integrar la justificaci\'on como fundamento te\'orico propio.
		\item Acr\'onimos \texttt{LBS}, \texttt{IoT}, \texttt{POI}, \texttt{GPS} definidos con \texttt{\textbackslash newcommand} pero \textbf{no usados} en el cuerpo. Eliminar o usarlos.
		\item El t\'itulo corto (``RutAI: AI-driven geolocated alarms\ldots'') no coincide con el t\'itulo principal; unificar.
		\item La Tabla~6 mezcla unidades heterog\'eneas (ms, mAh/h, bits, \%) sin columna expl\'icita de unidad, confuso para el lector.
	\end{itemize}
	
	\section{Veredicto}
	
	\begin{verdictbox}[Veredicto final]
		\textbf{Major revision.} El trabajo tiene m\'erito y se ha fortalecido considerablemente, pero no alcanza a\'un el umbral para JCR Q1/Q2 por: (a)~irreproducibilidad del banco de pruebas, (b)~validez discriminante no resuelta, (c)~ausencia de evaluaci\'on longitudinal en campo real y (d)~falta de evidencia visual del sistema.
	\end{verdictbox}
	
	% =============================================================
	\clearpage
	\part*{PARTE II --- Hoja de ruta del investigador senior}
	\addcontentsline{toc}{section}{PARTE II --- Hoja de ruta del investigador senior}
	% =============================================================
	
	A continuaci\'on expongo, con detalle operativo, \textbf{qu\'e datos faltan, d\'onde obtenerlos y con qu\'e protocolo}. El trabajo se organiza en cinco paquetes priorizados por impacto en la decisi\'on editorial.
	
	% =============================================================
	\section{Paquete A --- Rigor del banco de pruebas del bandido}
	\label{sec:paq-a}
	% =============================================================
	
	\emph{Prioridad: m\'axima.}
	
	\subsection{Construcci\'on del dataset semi-sint\'etico trazable}
	
	\textbf{Objetivo:} poder decir con precisi\'on de d\'onde viene cada curva de regret.
	
	\begin{actionbox}[A.1.1 --- Datasets p\'ublicos a utilizar]
		\begin{itemize}
			\item \emph{GeoLife GPS Trajectories} (Microsoft, 182 usuarios, 5 a\~nos): \\ \url{https://www.microsoft.com/en-us/download/details.aspx?id=52367}
			\item \emph{Taxi trajectories Porto / San Francisco Cabspotting}: disponible en CRAWDAD.
			\item \emph{Open Street Cam} / \emph{Mapillary} para condiciones de red.
			\item Para contexto latinoamericano: solicitar a la EPM de Quevedo o a la Alcald\'ia acceso a datos agregados de movilidad; alternativamente, \emph{Google Community Mobility Reports} (dominio p\'ublico).
		\end{itemize}
	\end{actionbox}
	
	\begin{actionbox}[A.1.2 --- Simulador determinista con c\'odigo publicado]
		Definir expl\'icitamente la funci\'on de recompensa del entorno simulado:
		\begin{lstlisting}[language=Python]
			r_t(a) = 0.7 * Bern(p_confirm(a | contexto)) + 0.3 * eta(a | contexto)
			p_confirm(a) = sigmoide(beta0 + beta1*tiempo_esperado + beta2*familiaridad)
			eta(a) = 1 - |tiempo_real - tiempo_planificado| / tiempo_planificado
		\end{lstlisting}
		Publicar en GitHub con \texttt{README}, \emph{tag} \texttt{v1.0.0} y DOI de Zenodo.
	\end{actionbox}
	
	\begin{actionbox}[A.1.3 --- Estad\'istica m\'as robusta]
		\begin{itemize}
			\item Incrementar r\'eplicas de 30 a \textbf{100 semillas} para estrechar los IC al 95\%.
			\item Reportar tambi\'en \textbf{regret por trip} (no s\'olo acumulado).
			\item A\~nadir an\'alisis de sensibilidad: variar $c\in\{0{,}5, 1, \sqrt{2}, 2\}$, $\lambda\in\{0, 0{,}1, 0{,}25, 0{,}5, 1{,}0\}$, horizonte de entrop\'ia $\in\{5, 10, 20, 50\}$. Presentar como \emph{heatmap} o superficie de respuesta.
			\item A\~nadir \textbf{contextual bandit LinUCB} (Li et al. 2010) como baseline adicional para demostrar que el UCB no contextual no ``juega en liga inferior''.
		\end{itemize}
	\end{actionbox}
	
	\textbf{Tiempo estimado:} 3--4 semanas de un investigador dedicado.
	
	\subsection{Test estad\'istico formal de la ventaja sobre Thompson sampling}
	
	El paper afirma ``$2{,}1\%$ mejor'' sin probar si es significativo. Protocolo sugerido:
	
	\begin{itemize}
		\item Test de Mann-Whitney U pareado por semilla sobre el regret final con $n=100$ r\'eplicas.
		\item Reportar tama\~no de efecto (\emph{Cliff's delta} o $r$).
		\item Si la diferencia no es significativa, \textbf{s\'e honesto y dilo}: ``UCB+entrop\'ia es estad\'isticamente indistinguible de Thompson en regret pero superior en entrop\'ia, que es el objetivo expl\'icito''.
	\end{itemize}
	
	% =============================================================
	\section{Paquete B --- Mediciones de sistema con rigor industrial}
	\label{sec:paq-b}
	% =============================================================
	
	\emph{Prioridad: cr\'itica.}
	
	\subsection{Latencia con condiciones documentadas}
	
	\begin{actionbox}[B.1 --- Protocolo de medici\'on de latencia]
		\textbf{Tres entornos de red}:
		\begin{itemize}
			\item LAN local (baseline ideal).
			\item 4G comercial (Claro/Movistar Ecuador o Vodafone Espa\~na, documentando \emph{throughput} medio).
			\item Wi-Fi dom\'estico saturado (descarga simult\'anea de 100\,MB como perturbaci\'on).
		\end{itemize}
		
		\textbf{Tres regiones del servidor}: Render US-East, Render Europa (si disponible) y VPS local (Contabo Quito $\sim$5\,ms). Esto aisla el componente de propagaci\'on.
		
		\textbf{Patr\'on de carga realista}: \texttt{k6} (preferido sobre Locust) con \emph{think-time} exponencial ($\lambda=1/3$\,s), arribo Poisson, rampa de $10\to 100$ usuarios en 60\,s, medir durante 5\,min de estado estacionario.
		
		\textbf{\emph{Warm-up} expl\'icito}: descartar los primeros 30\,s.
		
		\textbf{Reporte}: tabla con percentiles p50/p90/p95/p99 $\times$ 3 redes $\times$ 3 regiones, con IC \emph{bootstrap} al 95\%.
		
		\textbf{Herramientas}: \texttt{k6}, \texttt{tc netem} para simular latencia y p\'erdida controladas.
	\end{actionbox}
	
	\textbf{Tiempo estimado:} 1 semana.
	
	\subsection{Consumo energ\'etico con \texorpdfstring{$n\geq 5$}{n>=5} dispositivos}
	
	\begin{actionbox}[B.2 --- Protocolo de energ\'ia]
		\textbf{Selecci\'on de dispositivos} (m\'inimo 5, ideal 10):
		\begin{itemize}
			\item Gama baja: Samsung A04 / Motorola E13 / Xiaomi Redmi A2.
			\item Gama media: Samsung A34 / Pixel 6a / Motorola G73.
			\item Gama alta: Pixel 7/8, Samsung S23.
		\end{itemize}
		
		\textbf{Controles experimentales}:
		\begin{itemize}
			\item Brillo de pantalla al 50\%, \emph{timeout} de pantalla 30\,s.
			\item Sin app de fondo consumiendo red.
			\item Temperatura ambiente $22\pm 2$\,$^\circ$C.
			\item Bater\'ia al 80--85\% para evitar efectos de extremos.
		\end{itemize}
		
		\textbf{M\'etodo de medici\'on}:
		\begin{itemize}
			\item \texttt{adb shell dumpsys batterystats -{}-reset} al inicio.
			\item Sesi\'on de 2\,h con un gui\'on reproducible: crear recordatorio $\to$ caminar virtual (\texttt{adb shell geo fix}) $\to$ esperar trigger $\to$ repetir cada 15\,min.
			\item \texttt{adb shell dumpsys batterystats > session.log} al final.
			\item Battery Historian v2.1 para extraer mAh atribuibles a la app.
		\end{itemize}
		
		\textbf{Baselines adicionales}:
		\begin{itemize}
			\item FusedLocationProvider con \texttt{PRIORITY\_BALANCED\_POWER\_ACCURACY}.
			\item FusedLocationProvider con \texttt{PRIORITY\_HIGH\_ACCURACY}.
			\item Google Maps en \emph{background} compartiendo ubicaci\'on.
		\end{itemize}
		
		\textbf{An\'alisis estad\'istico}: ANOVA de medidas repetidas (condici\'on $\times$ dispositivo) + post-hoc Tukey HSD.
	\end{actionbox}
	
	\textbf{Tiempo estimado:} 2 semanas (ventana de medici\'on larga).
	
	\subsection{Geofencing con precisi\'on + recall + latencia del trigger}
	
	\begin{actionbox}[B.3 --- Protocolo de geofencing]
		\textbf{Conjunto de cercas}: 20 cercas en un barrio real (Quevedo, Granada, etc.), con radios 25/50/100\,m, georreferenciadas con DGPS (precisi\'on $<1$\,m) como \emph{ground truth}.
		
		\textbf{Perfiles de velocidad}:
		\begin{itemize}
			\item Peatonal lento (1,0\,m/s).
			\item Peatonal normal (1,4\,m/s).
			\item Bicicleta (4\,m/s).
			\item Veh\'iculo urbano (8,3\,m/s).
			\item Autopista (25\,m/s), opcional.
		\end{itemize}
		
		\textbf{Dise\~no}: por cada combinaci\'on cerca $\times$ radio $\times$ velocidad, m\'inimo \textbf{30 cruces reales} con dispositivo montado en veh\'iculo/bicicleta/persona, con \emph{log} sincronizado.
		
		\textbf{M\'etricas a reportar}:
		\begin{itemize}
			\item Precisi\'on, recall, F1.
			\item \textbf{Latencia del trigger}: $\Delta t$ entre cruce real y notificaci\'on.
			\item Percentiles de latencia p50/p95/p99.
		\end{itemize}
		
		\textbf{N objetivo}: m\'inimo 600 cruces (20 cercas $\times$ 3 radios $\times$ 4 velocidades $\times$ 30 repeticiones, reducible muestralmente si es inabarcable).
		
		\textbf{Equipo}: 1 asistente de campo durante 2 semanas; GPS de referencia Emlid Reach o similar ($\sim$500\,USD).
	\end{actionbox}
	
	% =============================================================
	\section{Paquete C --- Estudio TAM con potencia estad\'istica adecuada}
	\label{sec:paq-c}
	% =============================================================
	
	\emph{Prioridad: alta.}
	
	\subsection{Pre-registro del estudio}
	
	Antes de colectar, registrar el dise\~no en OSF (\url{https://osf.io}):
	
	\begin{itemize}
		\item Hip\'otesis primarias y secundarias.
		\item Criterios de inclusi\'on / exclusi\'on.
		\item Plan de an\'alisis estad\'istico.
		\item Umbrales de decisi\'on.
	\end{itemize}
	
	Esto es bien recibido por PMC y eleva la credibilidad del estudio.
	
	\subsection{Redise\~no del instrumento TAM}
	
	\textbf{Problema actual:} HTMT$>0{,}85$ indica que los \'items no discriminan constructos. Revisi\'on literal de cada \'item:
	
	\begin{enumerate}
		\item \textbf{Revisi\'on de contenido por panel de expertos} (m\'inimo 3 investigadores en HCI) con m\'etodo Lawshe (\emph{Content Validity Ratio}, CVR$\geq 0{,}62$ para $n=10$ jueces).
		\item \textbf{Refinamiento}: para cada par de constructos con HTMT$>0{,}85$, identificar \'items sem\'anticamente ambiguos y redactarlos para que separen conceptos. Ejemplos:
		\begin{itemize}
			\item PU original: ``La app es \'util para mis desplazamientos'' $\to$ comportamental-espec\'ifico: ``La app reduce el tiempo que dedico a planificar mis trayectos''.
			\item PS original: ``Me siento m\'as seguro usando la app'' $\to$ ``El monitoreo colaborativo me permite verificar si mis familiares llegaron a su destino''.
		\end{itemize}
		\item \textbf{Pilotaje con $n=20$}: s\'olo para validar comprensi\'on y tiempo de respuesta, no para inferencia.
	\end{enumerate}
	
	\subsection{Muestra definitiva con potencia suficiente}
	
	\textbf{C\'alculo de tama\~no}: para PLS-SEM con 5 constructos y rutas estructurales, m\'inimo $n=10\times$ n\'umero de rutas estructurales m\'as grande (regla de Hair). Con 4--5 rutas $\Rightarrow n\geq 50$. Para efecto peque\~no ($f^2=0{,}15$), potencia $0{,}80$, $\alpha=0{,}05\Rightarrow n\geq 138$ (G*Power~3.1).
	
	\textbf{Objetivo realista:} $n=150$ participantes, estratificados:
	
	\begin{itemize}
		\item 3 ciudades (Quevedo, Guayaquil y Quito; o Quevedo, Granada y una tercera).
		\item 50 por ciudad.
		\item Estratos por g\'enero (50/50), edad (18--29, 30--44, 45+), modalidad de transporte predominante.
	\end{itemize}
	
	\subsection{Protocolo de administraci\'on}
	
	\begin{enumerate}
		\item \textbf{Doble administraci\'on asincr\'onica} para reducir sesgo del investigador: primera vez presencial breve (15\,min) + 14 d\'ias de uso real + segunda administraci\'on del cuestionario en l\'inea.
		\item \textbf{Remuneraci\'on simb\'olica} (bono de datos m\'oviles, $\sim$5\,USD) para incentivar retenci\'on; declarar como COI.
		\item \textbf{Consentimiento escrito} con aprobaci\'on del Comit\'e de \'Etica de la universidad; n\'umero de aprobaci\'on mencionado en el \emph{paper}.
	\end{enumerate}
	
	\subsection{An\'alisis estad\'istico de nuevo nivel}
	
	\begin{itemize}
		\item PLS-SEM con SmartPLS~4 o \texttt{seminr} en R (c\'odigo abierto).
		\item Reportar: \emph{loadings}, AVE, CR, HTMT, VIF, $f^2$, $Q^2$, coeficientes de ruta con \emph{bootstrap} de 5000 muestras.
		\item IPMA (\emph{Importance-Performance Map Analysis}) para derivar implicaciones pr\'acticas.
		\item \textbf{MGA} (\emph{Multi-Group Analysis}) para comparar ciudades y grupos et\'areos.
	\end{itemize}
	
	\textbf{Tiempo estimado del paquete C completo:} 3--4 meses calendario.
	
	% =============================================================
	\section{Paquete D --- Evidencia de campo longitudinal}
	\label{sec:paq-d}
	% =============================================================
	
	\emph{Prioridad: alta. Sin este paquete, el t\'itulo ``Safer Urban Mobility'' es aspiracional, no emp\'irico.}
	
	\subsection{Estudio de campo de 6 semanas}
	
	\textbf{Dise\~no:} \emph{within-subjects} con fase de baseline y fase de intervenci\'on.
	
	\begin{itemize}
		\item \textbf{Semanas 1--2 (baseline)}: los participantes usan sus apps habituales (Google Maps, Waze). Se recolectan trazas GPS con consentimiento.
		\item \textbf{Semanas 3--6 (intervenci\'on)}: los mismos participantes usan \RUTAI{}.
	\end{itemize}
	
	\textbf{Muestra:} $n=50$ (potencia suficiente para test pareado con $d=0{,}5$ a $0{,}80$ de potencia).
	
	\subsection{Variables dependientes objetivas}
	
	\begin{enumerate}
		\item \textbf{Entrop\'ia de rutas real (no simulada)}: $H(U)$ calculada sobre las trayectorias GPS reales antes/despu\'es.
		\item \textbf{Distancia de Levenshtein} entre secuencias de rutas semanales (menor Levenshtein $=$ m\'as repetitiva).
		\item \textbf{\'Indice de regularidad temporal}: varianza de horarios de salida.
		\item \textbf{N\'umero de zonas \'unicas visitadas por semana}.
	\end{enumerate}
	
	\subsection{Variables dependientes subjetivas}
	
	\begin{enumerate}
		\item \textbf{State-Trait Anxiety Inventory (STAI)} aplicado semanalmente.
		\item \textbf{Perceived Stress Scale (PSS-10)}.
		\item \textbf{TAM} al final.
		\item \textbf{Diario de incidentes auto-reportados} (acoso, persecuci\'on percibida, robos).
	\end{enumerate}
	
	\subsection{An\'alisis estad\'istico}
	
	\begin{itemize}
		\item \emph{Paired-sample} $t$-test (o Wilcoxon si no es normal) para cada VD.
		\item ANCOVA controlando por covariables (g\'enero, edad, percepci\'on previa de inseguridad).
		\item \emph{Multilevel modeling} si el patr\'on var\'ia entre semanas.
	\end{itemize}
	
	\subsection{An\'alisis de seguridad real (si el presupuesto lo permite)}
	
	Si se encuentra colaboraci\'on con Polic\'ia Nacional u Observatorio de Seguridad Ciudadana, cruzar geolocalizaciones an\'onimas con \emph{hotspots} delictivos oficiales. M\'etrica: proporci\'on de tiempo en zonas de alto riesgo antes vs. despu\'es.
	
	\textbf{Tiempo:} 2 meses de trabajo de campo + 1 mes de an\'alisis.
	
	% =============================================================
	\section{Paquete E --- Materiales suplementarios y reproducibilidad}
	\label{sec:paq-e}
	% =============================================================
	
	\emph{Prioridad: alta.}
	
	\subsection{Repositorio p\'ublico}
	
	Estructura recomendada en GitHub:
	
	\begin{lstlisting}
		rutai/
		android-client/           (codigo Kotlin completo)
		backend/                  (FastAPI + requirements.txt con versiones fijas)
		simulator/                (banco de pruebas del bandido)
		analysis/                 (notebooks de analisis estadistico)
		data/
		tam-responses.csv       (anonimizado)
		technical-benchmark.csv
		geofencing-crossings.csv
		figures/                  (codigo para regenerar todas las figuras)
		docker-compose.yml        (un solo comando para levantar todo)
		LICENSE                   (MIT o Apache-2.0)
		CITATION.cff
		README.md
	\end{lstlisting}
	
	\textbf{DOI permanente:} Zenodo $\to$ \emph{tag} \texttt{v1.0.0} $\to$ DOI incrustado en el \emph{paper}.
	
	\subsection{Cuestionario TAM completo en ingl\'es y espa\~nol}
	
	PDF con los 25 \'items, su origen bibliogr\'afico, la escala, y la versi\'on traducida con \emph{back-translation} certificada.
	
	\subsection{Trazabilidad de requisitos}
	
	Matriz completa \emph{requisito $\leftrightarrow$ entrevista $\leftrightarrow$ fuente secundaria} como hoja Excel.
	
	\subsection{Capturas de pantalla de la app}
	
	M\'inimo cinco: \emph{login}, destinos guardados, selecci\'on de ruta con tres alternativas, mapa con geocerca activa, vista de grupo colaborativo con pings en tiempo real. Presentar en panel figura de $2\times 3$.
	
	\subsection{V\'ideo de demostraci\'on}
	
	V\'ideo de 3 minutos alojado en Zenodo / YouTube \emph{unlisted} que muestre las tres funcionalidades operando. Muchas revistas de PMC ya aceptan suplementos multimedia.
	
	% =============================================================
	\clearpage
	\section{Cronograma realista para alcanzar JCR Q1/Q2}
	% =============================================================
	
	\begin{longtable}{@{}p{3.6cm}p{2cm}p{4.5cm}p{3cm}@{}}
		\toprule
		\textbf{Paquete} & \textbf{Tiempo} & \textbf{Coste estimado} & \textbf{Prioridad} \\
		\midrule
		\endhead
		
		A. Banco de pruebas riguroso
		& 4 sem
		& Bajo (horas de investigador)
		& \textbf{Cr\'itica} \\
		
		B. Mediciones de sistema
		& 4 sem
		& Medio (\textasciitilde 500~USD GPS + 100~USD dispositivos)
		& \textbf{Cr\'itica} \\
		
		C. TAM $n=150$
		& 3--4 meses
		& Medio (\textasciitilde 750~USD participantes)
		& \textbf{Cr\'itica} \\
		
		D. Campo longitudinal 6 sem
		& 2 meses
		& Alto (\textasciitilde 2500~USD trabajo de campo)
		& \textbf{Alta} \\
		
		E. Reproducibilidad
		& 2 sem
		& Bajo
		& \textbf{Alta} \\
		
		\bottomrule
	\end{longtable}
	
	\textbf{Calendario sugerido:}
	
	\begin{itemize}
		\item Empezar paquetes A y B en paralelo durante las primeras 4 semanas.
		\item Iniciar C al final de la semana 2 (solapamiento).
		\item Iniciar D al finalizar C.
		\item Entrega del manuscrito revisado: \textbf{6 meses desde el inicio del plan}.
	\end{itemize}
	
	% =============================================================
	\section{Estrategia alternativa si los recursos son limitados}
	% =============================================================
	
	Si el equipo no dispone de presupuesto para los paquetes C y D completos, una alternativa viable ser\'ia:
	
	\begin{enumerate}
		\item \textbf{Fortalecer A y B} (mediciones t\'ecnicas) para dar peso al lado ingenier\'ia.
		\item \textbf{Mantener el TAM como exploratorio} ($n=25$) pero cambiar el \emph{framing}: el manuscrito pasa a ser un \emph{system paper} con mediciones t\'ecnicas s\'olidas $+$ evaluaci\'on cualitativa preliminar.
		\item \textbf{Redirigir la revista objetivo} a una donde el aspecto t\'ecnico pese m\'as que el humano:
		\begin{itemize}
			\item \emph{Future Generation Computer Systems} (Elsevier, Q1, IF~7.5): reforzando el \'angulo de sistemas distribuidos.
			\item \emph{Journal of Network and Computer Applications} (Elsevier, Q1, IF~7.7): reforzando el \'angulo de arquitectura m\'ovil.
			\item \emph{Sensors} (MDPI, Q2, IF~3.4): m\'as tolerante con $n$ peque\~no, \emph{open access}.
			\item \emph{Electronics} (MDPI, Q2, IF~2.6).
		\end{itemize}
	\end{enumerate}
	
	Esta ruta es m\'as r\'apida (3 meses) pero tiene menor impacto que PMC Q2.
	
	% =============================================================
	\section{Conclusi\'on del investigador senior}
	% =============================================================
	
	\RUTAI{} tiene los ingredientes de un buen \emph{paper} Q2. Lo que le falta no es creatividad ni ambici\'on, sino \textbf{evidencia cuantitativa y reproducibilidad a la altura del est\'andar actual}. Con los paquetes A, B, C y E bien ejecutados, el manuscrito puede aspirar razonablemente a PMC en segunda ronda. Si se a\~nade el paquete D, podr\'ia incluso escalar a revista Q1 (\emph{FGCS} o \emph{Computers \& Security}). Sin esos datos, el riesgo de rechazo en Q2 supera el $60\%$.
	
	El orden importa: \textbf{primero A+B (rigurosas pero ejecutables en 4 semanas), despu\'es C en paralelo con E, y D si los recursos lo permiten}. Evitar la tentaci\'on de enviar con los datos actuales y esperar que el revisor sea benevolente; los revisores JCR de 2026 no lo son.
	
	\begin{verdictbox}[Recomendaci\'on final al equipo autor]
		\begin{enumerate}
			\item No enviar el manuscrito en su estado actual.
			\item Ejecutar los paquetes A y B en paralelo durante las primeras 4 semanas.
			\item Dise\~nar y pre-registrar el paquete C mientras A y B se ejecutan.
			\item Desplegar E (repositorio, capturas, v\'ideo) en la \'ultima semana antes del env\'io.
			\item Si existe presupuesto, a\~nadir D y apuntar a una revista Q1.
		\end{enumerate}
	\end{verdictbox}
	
\end{document}